\section{Bölüm sonu çözümlü problemler}

\begin{enumerate}
  \item \textbf{Soru--yöntem eşleştirmesi.} Aynı öğrencilerin ön-test ve son-test puanları, iki bağımsız program grubunun son-test puanları ve çalışma saatiyle başarı puanı incelenecektir. Üç temel yöntemi seçin.

  \textbf{Çözüm:} Aynı öğrencilerin değişimi için eşleştirilmiş $t$ testi; bağımsız grupların ortalama farkı için Welch $t$ testi; iki nicel değişkenin doğrusal ilişkisi için saçılım grafiğiyle birlikte Pearson korelasyonu veya basit doğrusal regresyon kullanılır.

  \item \textbf{Tutarsız görünen sonuçlar.} Grup farkı için $d=0{,}60$, yüzde 95 güven aralığı $[-0{,}2;1{,}4]$ ve $p=0{,}14$ bulunmuştur. Sonucu bütünleştirin.

  \textbf{Çözüm:} Nokta tahmini orta büyüklükte görünse de aralık küçük ters etkiden büyük pozitif etkiye kadar geniştir ve sıfırı içerir. $p$-değeri seçilen eşikte kanıtın sınırlı olduğunu gösterir. Sonuç ``etki yoktur'' değil, tahmin hassasiyetinin yetersiz olduğu biçiminde sunulmalıdır.

  \item \textbf{Tekrarlanabilir analiz paketi.} Ham veri üzerinde iki elle düzeltme yapılmış, fakat değişiklikler kaydedilmemiştir. Süreci yeniden üretilebilir hâle getirmek için ne yapılmalıdır?

  \textbf{Çözüm:} Ham dosya değişmeden korunmalı; düzeltmeler doğrulama kuralları ve gerekçeleriyle kodlanarak temiz veri yeniden üretilmelidir. Veri sözlüğü, analiz günlüğü, yazılım sürümleri, kod, rassal tohumlar ve çıktı üretim adımları aynı proje yapısında saklanmalıdır.
\end{enumerate}